\documentclass[11pt,a4paper]{article}
\usepackage[utf8]{inputenc}
\usepackage[T1]{fontenc}
\usepackage{lmodern,amsmath,amssymb,amsthm,mathtools,booktabs,hyperref,enumitem,microtype}
\usepackage[margin=1in]{geometry}
\usepackage[numbers,sort&compress]{natbib}

\newtheorem{definition}{Definition}[section]
\newtheorem{proposition}[definition]{Proposition}
\theoremstyle{remark}
\newtheorem{remark}[definition]{Remark}

\newcommand{\Frame}{\mathcal F}
\newcommand{\Seldon}{\mathcal S}

\title{Seldon Manifolds --- Psychological Grounding\\\large Forward Models, Acceptable Regions, and the Aperture of One Embodiment}
\author{Matthew Parslow\\Independent Researcher}
\date{\textbf{Draft \textemdash\ draft-2026-08-29.1}\\[2pt]29 August 2026 \textbullet\ document revision 2}

\usepackage{fancyhdr}
\pagestyle{fancy}
\fancyhf{}
\fancyfoot[C]{\thepage}
\fancyfoot[R]{\scriptsize\textsc{draft-2026-08-29.1\,r2}}
\renewcommand{\headrulewidth}{0pt}
\fancypagestyle{plain}{\fancyhf{}\fancyfoot[C]{\thepage}%
  \fancyfoot[R]{\scriptsize\textsc{draft-2026-08-29.1\,r2}}\renewcommand{\headrulewidth}{0pt}}
\begin{document}
\maketitle
\begin{center}
\fbox{\parbox{0.92\textwidth}{\small
\textbf{DRAFT --- draft-2026-08-29.1 r2} (29 August 2026).
Circulated for comment, not as a finished result. This is the grounding companion to
\emph{Seldon Manifolds --- Theory}, which sets out what a Seldon manifold is and how it works;
this paper asks only whether one particular substrate implements anything of that shape. Sections
marked as falsifiability commitments are the intended points of attack.}}
\end{center}
\vspace{1em}

\begin{abstract}
The Seldon manifold is a construct built to do forecasting, not a model of human cognition, so
nothing here bears on whether it works. What this companion asks is narrower and worth asking on
its own account: does a biological nervous system implement anything of that shape, and if so with
what horizon, what resolution, and over which dimensions. The answer is that several components are
implemented and well measured --- a forward model, an acceptable \emph{region} rather than a target
with correction only on deviations that would leave it, resolution degrading with temporal distance
in a known functional form, and boundaries indexed to the body rather than to the world --- while
the composite is not. The measured motor horizon is on the order of a hundred and fifty
milliseconds and exists to cancel sensory delay; deliberative search runs about three steps deep and
is truncated by a reflex rather than by a cost calculation; and within-frame forecasting is good
where a domain has recurrent structure to hold and worth nothing where it does not. The honest
summary is that the pieces are real, the human instantiation is narrow, and its quality is
frame-relative rather than general. One line of evidence the construct would most naturally have
leaned on --- that representations coarsen with psychological distance --- does not survive its
best-powered test, and is not relied upon here.
\end{abstract}

\section{Scope}\label{sec:scope}
This paper is a grounding and not a validation, and the distinction is load-bearing. The construct
it grounds was built as engineering, to forecast; it makes no claim about how people think, so
evidence that people forecast badly would not count against it. What such evidence does is
characterise one implementation. Where the psychology is unflattering it is therefore reported as a
measurement of that implementation's limits rather than as an objection, and where it is
flattering it is reported without the suggestion that the construct needed it.

Two consequences follow for how this paper should be read. Findings are organised by the component
they bear on --- horizon, resolution, represented dimensions, fields, the acceptable region, and the
capacity to reshape rather than navigate --- rather than by the literature they come from, so several
mature research traditions appear here in pieces. And a component may be well evidenced while the
composite is not; nothing below argues that any nervous system holds a Seldon manifold, only that
each part of one has an implemented counterpart somewhere, mostly in different systems and at wildly
different timescales.

\section{Forward models, and a horizon of milliseconds}\label{sec:motor}
The least controversial forward model in biology is the motor one. A forward model in this sense
takes the current state and the outgoing motor command and predicts the next state
\citep{miallwolpert1996}, and the evidence that the nervous system runs one is behavioural rather
than inferential: sensory consequences of self-produced movement are attenuated in proportion to
how well they are predicted, so that perturbing the timing or trajectory of a self-administered
tactile stimulus restores the sensation it would otherwise have suppressed
\citep{blakemore1999}, with a delay of roughly a hundred milliseconds sufficient to break the
attenuation \citep{kilteni2019}. Prediction is also learned before control is
\citep{flanagan2003}, which separates the forward model from the controller that uses it.

The horizon of that model has been measured, and it is short. Disrupting the lateral cerebellum by
stimulation makes reaches behave as though planned from a hand position around $138\,$ms out of
date, replicated at $134\,$ms, with the authors' upper estimate $150$--$160\,$ms
\citep{miall2007}. That figure is not an incidental limitation but the model's reason for
existing: the feedback loops it compensates run at twenty to sixty milliseconds spinally and around
a hundred visually, and a predictor that cancels them needs exactly that much lead. The motor
forward model is a delay compensator rather than a simulator of futures, and the gap between it and
a deliberative horizon is not bridged by any evidence in this paper.

\begin{remark}[What this does and does not license]\label{rem:motor-limits}
It licenses the claim that a nervous system maintains a predictive model of its own next states and
acts on the prediction rather than on the measurement. It does not license any claim about the
length of a planning horizon, and a reader who takes the two together is making an architectural
argument that the evidence here does not supply.
\end{remark}

\section{An acceptable region, and correction only on departure}\label{sec:region}
The theory paper holds that a frame may represent a region of acceptable futures and intervene only
when its projection leaves that region. Motor control implements something with this shape, and the
implementation is measured twice over.

The first is structural. Partitioning the variability of a practised movement into components that
do and do not affect the task outcome shows that variance within the task-equivalent subspace is
left alone while variance that would change the outcome is corrected --- the
\emph{uncontrolled manifold} \citep{scholz1999}. The nervous system is not stabilising a
trajectory; it is stabilising a region and permitting drift inside it.

The second is theoretical and derived rather than assumed. Under optimal feedback control the
resulting policy exhibits what its authors name the \emph{minimal intervention} principle:
deviations from the average trajectory are corrected only when they interfere with task
performance \citep{todorov2002}.

\begin{remark}[Three cautions on the correspondence]\label{rem:mip-cautions}
The principle is \emph{minimal} intervention, not minimum: its criterion is task relevance, not
magnitude, and the same analysis permits deviations to be \emph{magnified} where doing so reduces
cost. It is therefore not a smallest-sufficient-push result and should not be cited as one. And the
controller does not project forward at decision time: the horizon is consumed offline in a backward
pass, leaving an online policy that is a single operation on a state estimate. An architecture that
re-projects at each step is closer to receding-horizon control, which \citet{todorov2002} treat as a
separate approach.
\end{remark}

The idea that acting well means keeping a projection inside a region is older than either result.
Driving was described in 1938 in terms of a \emph{field of safe travel}, an explicitly spatial
region of acceptable trajectories within which the driver steers \citep{gibsoncrooks1938}, a
formulation whose computational descendants still maintain a perceived risk field below a threshold
\citep{kolekar2020}.

\section{The acceptable region is Simon's claim, not a simplification of it}\label{sec:satisficing}
Satisficing is usually reported as accepting the first option above a threshold, which is a scalar
cutoff on one dimension. That is not what the original says. The aspiration level adapts --- the
balance between what is needed and what is available is maintained by \emph{raising and lowering}
aspiration levels --- and the criterion is conjunctive across needs, a path being sought that permits
satisfaction at some specified level of \emph{all} of them \citep{simon1956}. Simultaneous minima on
every need dimension describe a region in a multi-dimensional space, and the fixed scalar version is
explicitly the disposable simplification of a toy organism.

The formal development followed: aspiration adaptation theory treats goals as a \emph{vector} with an
aspiration on each component, adjusted against feasibility \citep{selten1998}, and non-compensatory
conjunctive models --- a cutoff per attribute with no trade-offs between them --- have been fitted to
choice process data in preference to compensatory alternatives \citep{einhorn1970,stuttgen2012}.
That the aspirations adapt to the environment rather than sitting fixed is itself measured
\citep{caplin2011}.

\begin{remark}[Where the region is genuinely absent]\label{rem:thresholds}
Several familiar constructs look like acceptable regions and are not. A reference point is one point
on one axis with a kink, below which outcomes are treated continuously differently rather than
categorically \citep{kahneman1979}; safety-first criteria minimise the probability of falling below a
level, which is a half-line \citep{roy1952}; and aspiration in the risk-preference literature, while
a genuine second criterion rather than a reshaped utility function, remains a scalar cutoff
\citep{lopes1999}. The distinction between a threshold on a scalar and a region in several
dimensions is exactly where an account can go wrong: the best-known attempt to explain accident
rates by a single target level of risk failed, and failed because it collapsed a field to one scalar
on one variable \citep{wilde1982,evans1986}.
\end{remark}

\section{Resolution falls with horizon}\label{sec:resolution}
The theory paper treats resolution and horizon as separate axes under a shared budget, with
resolution characteristically degrading as horizon extends. There is a measured version of this with
a functional form. In interval timing the standard deviation of an estimate is proportional to the
interval being estimated, so the coefficient of variation stays constant as the interval grows
\citep{gibbon1977,rakitin1998,zeki2019}. Absolute uncertainty therefore scales with distance: the
further out, the coarser, in a fixed ratio. The result is fifty years old, holds across species, and
is not in dispute.

\begin{remark}[A stronger candidate that does not survive]\label{rem:clt}
There is an obvious alternative anchor, and it should be named because a reader will otherwise
supply it. Construal level theory holds that the further an object is from direct experience the
more abstract its representation \citep{tropeliberman2010}, which is resolution falling with
distance stated for cognition generally rather than for timing. A meta-analysis of over two hundred
experiments reported a reliable medium effect \citep{soderberg2015}. But a registered multi-laboratory
replication across four distance types, with nearly twelve thousand participants, found effects of
$d\approx0.03$--$0.08$ with the social effect reversed, and a bias-corrected re-analysis of the
existing literature finds strong evidence of publication bias \citep{calderon2025,maier2022}. The
theory's authors contest the replication \citep{liberman2026}, and both critical papers are
preprints, so the accurate statement is that the effect is in serious dispute rather than that it is
refuted --- and either way it cannot carry weight. The timing result above is narrower and secure,
and is used instead.
\end{remark}

\section{Horizon is plural and frame-relative}\label{sec:horizon}
A single agent does not have one horizon. Cortical regions differ systematically in the window over
which they integrate information, from tens of milliseconds in sensory areas to many seconds in
higher-order ones \citep{hasson2008,hasson2015}, and predictive representations differ in reach
along the same gradient, with anterior prefrontal cortex showing the largest predictive horizons and
posterior hippocampus the smallest \citep{brunec2022}. Different subsystems hold different horizons
over the same world, which is the theory paper's frame-relativity appearing within one skull rather
than between agents.

\section{The manifold is indexed to the body}\label{sec:body}
The strongest support in this paper is for the claim that different agents hold different manifolds
of one situation, with no frame-free manifold that theirs approximate. The boundary between a
stairway that can be climbed and one that cannot is not a property of the stairway: it scales to the
observer's leg length and matches biomechanical prediction \citep{warren1984}, and the boundary
between an aperture that can be walked through and one requiring rotation sits at a constant ratio
of aperture to shoulder width, recovered from body-scaled optical information \citep{warrenwhang1987}.
The same physical situation has measurably different boundaries for different bodies, and there is no
body-free version of the boundary for either of them to be wrong about.

\begin{remark}[The extension is not free]\label{rem:affordance-limit}
This is the perception of immediately available action, not a multi-step future. That the same
frame-relativity holds of extended horizons is an argument the theory paper makes and this
literature does not supply.
\end{remark}

\section{Bounded search, and how it stops}\label{sec:search}
Measured planning depth is shallow. Where it has been fitted directly it runs to roughly three steps,
and subjects trade depth against precision under a near-constant workload, preferring to reduce
depth rather than compute less precisely \citep{snider2015}; deep problems are approached by
concatenating shallow searches rather than by deepening them, with subgoal selection chosen so as to
nearly minimise the cost of computing values \citep{huys2015}.

How the search stops is less comfortable for any account that wants a cost-benefit rule. Branches are
pruned reflexively on encountering a large loss, and the pruning persists when it is made
overwhelmingly counterproductive \citep{huys2012}. The same subjects show near-optimal subgoal
selection and a hard-wired affective guillotine in the same task \citep{huys2015}. The honest reading
is that stopping is heuristic in mechanism and only approximately rational in aggregate.

\begin{remark}[Shallow search is not a shallow forward model]\label{rem:sr}
It would be a mistake to infer from three-step search that the forward representation is
impoverished, and there is a better-supported alternative. People appear to hold a precomputed,
horizon-discounted predictive map of expected future occupancy --- flexible when rewards change,
insensitive when the sequence of states changes --- rather than either searching a tree or caching a
policy \citep{momennejad2017}. Such a map is much closer to a manifold with fields over it than tree
search is, and it explains shallow measured search without impoverishing the representation. The
caveat is that it is a compressed occupancy prediction rather than a branching structure with
explicit alternatives, so it under-delivers on reachable futures in the plural.
\end{remark}

\section{Within-frame quality, and why it does not transfer}\label{sec:expertise}
Forecasting quality inside a frame can be very high, and it is not bought by searching further.
Replicating the classical result across skill levels finds that stronger players search
\emph{faster}, not deeper, with pattern recognition rather than search the principal discriminator
\citep{degroot1965,connors2011}; the sharpest demonstration is that a grandmaster's rated strength
falls only slightly when playing six boards simultaneously against tournament-length thinking time
\citep{gobetsimon1996ps}. Depth differences with skill are not zero \citep{campitelli2004}, so
recognition should be called the principal discriminator rather than the only one.

The anticipation is genuinely prospective where it has been isolated: skilled performers respond to
advance kinematic information before the event they are anticipating occurs, with the effect
established across dozens of studies and hundreds of effect sizes
\citep{muller2006,mann2007}, and the advantage disappears when attention is forced elsewhere
\citep{muller2010}.

It also does not transfer. Far transfer from trained domains shrinks as study design improves and
vanishes against active controls \citep{sala2017cdps,sala2017lb,melbylervag2016}. The mechanism
predicts this rather than merely tolerating it: a store of domain-structured patterns has nothing to
decode a foreign domain against. The complementary case is forecasting in a domain without stable
recurrent structure, where experts do no better than simple statistical models or informed crowds
\citep{forecasting2023} --- though within the same literature domain expertise still predicts
accuracy, and calibration improves with training and aggregation \citep{mellers2014}. Where
structure exists the map forms; where it does not, expertise buys nothing.

\begin{remark}[One control is weaker than usually reported]\label{rem:random-boards}
The random-position control is often cited as showing that expert advantage vanishes without domain
structure. It does not fully vanish: strong players retain a smaller but real advantage on random
material, and it generalises across fields \citep{gobetsimon1996pbr,sala2017mc}. The non-transfer
claim above rests on the far-transfer meta-analyses, not on this control.
\end{remark}

\section{Blind spots, reshaping, and the growth of the dimension set}\label{sec:other}
Three further components have implemented counterparts.

\emph{Blind spots are real and are not introspectively accessible.} People report understanding
complex phenomena with far greater precision and depth than they possess, and discover this only on
attempting explanation \citep{rozenblit2002}; attended observers miss salient unexpected events
entirely \citep{simons1999}; and adults who can distinguish their own knowledge from another's
routinely fail to deploy that distinction online \citep{keysar2003}. This supports the theory
paper's treatment of a blind spot as a distinct deficit from coarseness, and rules out any mechanism
in which a frame audits its own dimension set from inside.

\emph{Reshaping rather than navigating is implemented.} Committing in advance so as to remove a
future option is a documented and modelled behaviour \citep{rachlin1972,ainslie1975}, and the
broader category of changing one's situation rather than choosing within it is treated as a distinct
class of self-control \citep{duckworth2016}.

\emph{The dimension set can grow.} Perceptual learning includes attentional weighting, differentiation
and unitisation, and people demonstrably learn to separate dimensions that were previously fused
\citep{goldstone1998,goldstone2001}. A represented dimension set is therefore not fixed, which the
theory paper requires and does not otherwise evidence.

\section{What the evidence does not support}\label{sec:limits}
Four things should be stated plainly, since a grounding paper that reports only the matches is not
worth reading.

\emph{The composite is not evidenced.} Every component above is implemented somewhere, but in
different systems, at timescales separated by three orders of magnitude, and no study assembles them.
That a nervous system holds a Seldon manifold is not shown here and is not claimed.

\emph{Whether the machinery is required is unresolved.} Patients with hippocampal amnesia have been
reported both as unable to construct coherent imagined experiences \citep{hassabis2007} and as
unimpaired at imagining the future \citep{squire2010}, with a third position arguing that the verbal
measures used cannot adjudicate between them \citep{miloyan2019}. The dispute is live. It is
compatible with the theory paper's treatment of the manifold as optional machinery, but it should be
presented as unresolved rather than as settled favourably.

\emph{Forecasting error is on magnitude, not structure --- and this is a reframing.} People
systematically overestimate the intensity and duration of their future emotional reactions
\citep{gilbert1998,wilsongilbert2005}. Much of the effect appears to be an artefact of how the
question is interpreted \citep{levine2012}, and when measured in daily life people predict correctly
\emph{which} events will make them feel better or worse while erring on absolute levels
\citep{moeck2026}. A manifold whose shape is right and whose field magnitudes are miscalibrated is
what this describes. That is a legitimate reading and it is still a reading, offered as such.

\emph{The per-agent axes are not established.} The theory paper allows different frames to differ
independently in horizon, resolution and represented dimensions. No study assembled here measures
those in one sample, and at least one natural axis does not separate --- imagery vividness and
richness of future simulation covary rather than dissociating. Separable per-agent axes are a
synthesis, not a finding.

\section{Falsifiability commitments}\label{sec:falsifiability}
\begin{itemize}[leftmargin=1.4em]
\item \textbf{A predictive forward model.} Fails if the attenuation of self-produced sensation does
not track prediction error, or if prediction is not separable from control.
\item \textbf{Region rather than target.} Fails if task-irrelevant variability is corrected as
strongly as task-relevant variability, i.e.\ if there is no uncontrolled subspace.
\item \textbf{Resolution scaling with distance.} Fails if the coefficient of variation of temporal
estimates is not approximately constant across intervals.
\item \textbf{Body-indexed boundaries.} Fails if perceived action boundaries do not scale with the
relevant bodily dimension, or if one boundary fits all bodies.
\item \textbf{Structure-dependent forecasting.} Fails if within-domain forecasting advantage is
independent of the domain's recurrent structure, or if it transfers to unrelated domains under
active controls.
\item \textbf{Scope.} This paper is falsified as over-claiming if any section is read as evidence
that a nervous system holds a Seldon manifold, or as validation of a construct that does not require
it.
\end{itemize}

\begin{thebibliography}{99}
\bibitem[Ainslie(1975)]{ainslie1975}
George Ainslie.
\newblock Specious reward: a behavioral theory of impulsiveness and impulse control.
\newblock \emph{Psychological Bulletin}, 82(4):463--496, 1975.

\bibitem[Blakemore et al.(1999)]{blakemore1999}
Sarah-Jayne Blakemore, Chris D. Frith, and Daniel M. Wolpert.
\newblock Spatio-temporal prediction modulates the perception of self-produced stimuli.
\newblock \emph{Journal of Cognitive Neuroscience}, 11(5):551--559, 1999.

\bibitem[Brunec and Momennejad(2022)]{brunec2022}
Iva K. Brunec and Ida Momennejad.
\newblock Predictive representations in hippocampal and prefrontal hierarchies.
\newblock \emph{Journal of Neuroscience}, 42(2):299--312, 2022.

\bibitem[Calderon et al.(2025)]{calderon2025}
Sofia Calderon, Erik Mac Giolla, Karl Ask, and Timothy J. Luke.
\newblock Effects of psychological distance on mental abstraction: a registered report of four tests of construal level theory.
\newblock PsyArXiv preprint, 2025.

\bibitem[Campitelli and Gobet(2004)]{campitelli2004}
Guillermo Campitelli and Fernand Gobet.
\newblock Adaptive expert decision making: skilled chess players search more and deeper.
\newblock \emph{ICGA Journal}, 27(4):209--216, 2004.

\bibitem[Caplin et al.(2011)]{caplin2011}
Andrew Caplin, Mark Dean, and Daniel Martin.
\newblock Search and satisficing.
\newblock \emph{American Economic Review}, 101(7):2899--2922, 2011.

\bibitem[Connors et al.(2011)]{connors2011}
Michael H. Connors, Bruce D. Burns, and Guillermo Campitelli.
\newblock Expertise in complex decision making: the role of search in chess 70 years after de Groot.
\newblock \emph{Cognitive Science}, 35(8):1567--1579, 2011.

\bibitem[de Groot(1965)]{degroot1965}
Adriaan D. de Groot.
\newblock \emph{Thought and Choice in Chess}.
\newblock Mouton, The Hague, 1965. Originally published in Dutch, 1946.

\bibitem[Duckworth et al.(2016)]{duckworth2016}
Angela L. Duckworth, Tamar Szab\'o Gendler, and James J. Gross.
\newblock Situational strategies for self-control.
\newblock \emph{Perspectives on Psychological Science}, 11(1):35--55, 2016.

\bibitem[Einhorn(1970)]{einhorn1970}
Hillel J. Einhorn.
\newblock The use of nonlinear, noncompensatory models in decision making.
\newblock \emph{Psychological Bulletin}, 73(3):221--230, 1970.

\bibitem[Evans(1986)]{evans1986}
Leonard Evans.
\newblock Risk homeostasis theory and traffic accident data.
\newblock \emph{Risk Analysis}, 6(1):81--94, 1986.

\bibitem[Flanagan et al.(2003)]{flanagan2003}
J. Randall Flanagan, Philipp Vetter, Roland S. Johansson, and Daniel M. Wolpert.
\newblock Prediction precedes control in motor learning.
\newblock \emph{Current Biology}, 13(2):146--150, 2003.

\bibitem[Forecasting Collaborative(2023)]{forecasting2023}
Forecasting Collaborative.
\newblock Insights into the accuracy of social scientists' forecasts of societal change.
\newblock \emph{Nature Human Behaviour}, 7(4):484--501, 2023.

\bibitem[Gibbon(1977)]{gibbon1977}
John Gibbon.
\newblock Scalar expectancy theory and Weber's law in animal timing.
\newblock \emph{Psychological Review}, 84(3):279--325, 1977.

\bibitem[Gibson and Crooks(1938)]{gibsoncrooks1938}
James J. Gibson and Laurence E. Crooks.
\newblock A theoretical field-analysis of automobile-driving.
\newblock \emph{American Journal of Psychology}, 51(3):453--471, 1938.

\bibitem[Gilbert et al.(1998)]{gilbert1998}
Daniel T. Gilbert, Elizabeth C. Pinel, Timothy D. Wilson, Stephen J. Blumberg, and Thalia P. Wheatley.
\newblock Immune neglect: a source of durability bias in affective forecasting.
\newblock \emph{Journal of Personality and Social Psychology}, 75(3):617--638, 1998.

\bibitem[Gobet and Simon(1996a)]{gobetsimon1996ps}
Fernand Gobet and Herbert A. Simon.
\newblock The roles of recognition processes and look-ahead search in time-constrained expert problem solving.
\newblock \emph{Psychological Science}, 7(1):52--55, 1996.

\bibitem[Gobet and Simon(1996b)]{gobetsimon1996pbr}
Fernand Gobet and Herbert A. Simon.
\newblock Recall of random and distorted chess positions: implications for the theory of expertise.
\newblock \emph{Psychonomic Bulletin \& Review}, 3(2):159--163, 1996.

\bibitem[Goldstone(1998)]{goldstone1998}
Robert L. Goldstone.
\newblock Perceptual learning.
\newblock \emph{Annual Review of Psychology}, 49:585--612, 1998.

\bibitem[Goldstone and Steyvers(2001)]{goldstone2001}
Robert L. Goldstone and Mark Steyvers.
\newblock The sensitization and differentiation of dimensions during category learning.
\newblock \emph{Journal of Experimental Psychology: General}, 130(1):116--139, 2001.

\bibitem[Hassabis et al.(2007)]{hassabis2007}
Demis Hassabis, Dharshan Kumaran, Seralynne D. Vann, and Eleanor A. Maguire.
\newblock Patients with hippocampal amnesia cannot imagine new experiences.
\newblock \emph{Proceedings of the National Academy of Sciences}, 104(5):1726--1731, 2007.

\bibitem[Hasson et al.(2008)]{hasson2008}
Uri Hasson, Eunice Yang, Ignacio Vallines, David J. Heeger, and Nava Rubin.
\newblock A hierarchy of temporal receptive windows in human cortex.
\newblock \emph{Journal of Neuroscience}, 28(10):2539--2550, 2008.

\bibitem[Hasson et al.(2015)]{hasson2015}
Uri Hasson, Janice Chen, and Christopher J. Honey.
\newblock Hierarchical process memory: memory as an integral component of information processing.
\newblock \emph{Trends in Cognitive Sciences}, 19(6):304--313, 2015.

\bibitem[Huys et al.(2012)]{huys2012}
Quentin J. M. Huys, Neir Eshel, Elizabeth O'Nions, Luke Sheridan, Peter Dayan, and Jonathan P. Roiser.
\newblock Bonsai trees in your head: how the Pavlovian system sculpts goal-directed choices by pruning decision trees.
\newblock \emph{PLoS Computational Biology}, 8(3):e1002410, 2012.

\bibitem[Huys et al.(2015)]{huys2015}
Quentin J. M. Huys, Ni\'all Lally, Paul Faulkner, Neir Eshel, Erich Seifritz, Samuel J. Gershman, Peter Dayan, and Jonathan P. Roiser.
\newblock Interplay of approximate planning strategies.
\newblock \emph{Proceedings of the National Academy of Sciences}, 112(10):3098--3103, 2015.

\bibitem[Kahneman and Tversky(1979)]{kahneman1979}
Daniel Kahneman and Amos Tversky.
\newblock Prospect theory: an analysis of decision under risk.
\newblock \emph{Econometrica}, 47(2):263--291, 1979.

\bibitem[Keysar et al.(2003)]{keysar2003}
Boaz Keysar, Shuhong Lin, and Dale J. Barr.
\newblock Limits on theory of mind use in adults.
\newblock \emph{Cognition}, 89(1):25--41, 2003.

\bibitem[Kilteni et al.(2019)]{kilteni2019}
Konstantina Kilteni, Christian Houborg, and H. Henrik Ehrsson.
\newblock Rapid learning and unlearning of predicted sensory delays in self-generated touch.
\newblock \emph{eLife}, 8:e42888, 2019.

\bibitem[Kolekar et al.(2020)]{kolekar2020}
Sarvesh Kolekar, Joost de Winter, and David Abbink.
\newblock Human-like driving behaviour emerges from a risk-based driver model.
\newblock \emph{Nature Communications}, 11:4850, 2020.

\bibitem[Levine et al.(2012)]{levine2012}
Linda J. Levine, Heather C. Lench, Robin L. Kaplan, and Martin A. Safer.
\newblock Accuracy and artifact: reexamining the intensity bias in affective forecasting.
\newblock \emph{Journal of Personality and Social Psychology}, 103(4):584--605, 2012.

\bibitem[Liberman and Levit-Mor(2026)]{liberman2026}
Nira Liberman and Or Levit-Mor.
\newblock The effect of psychological distance on level of construal: what do we learn from the CLIMR failed replication?
\newblock PsyArXiv preprint, 2026.

\bibitem[Lopes and Oden(1999)]{lopes1999}
Lola L. Lopes and Gregg C. Oden.
\newblock The role of aspiration level in risky choice.
\newblock \emph{Journal of Mathematical Psychology}, 43(2):286--313, 1999.

\bibitem[Maier et al.(2022)]{maier2022}
Maximilian Maier, Franti\v{s}ek Barto\v{s}, Minjeong Oh, Eric-Jan Wagenmakers, David Shanks, and Adam J. L. Harris.
\newblock Adjusting for publication bias reveals that evidence for and size of construal level theory effects is substantially overestimated.
\newblock PsyArXiv preprint, 2022.

\bibitem[Mann et al.(2007)]{mann2007}
Derek T. Y. Mann, A. Mark Williams, Paul Ward, and Christopher M. Janelle.
\newblock Perceptual-cognitive expertise in sport: a meta-analysis.
\newblock \emph{Journal of Sport and Exercise Psychology}, 29(4):457--478, 2007.

\bibitem[Melby-Lerv\aa g et al.(2016)]{melbylervag2016}
Monica Melby-Lerv\aa g, Thomas S. Redick, and Charles Hulme.
\newblock Working memory training does not improve performance on measures of intelligence or other measures of ``far transfer''.
\newblock \emph{Perspectives on Psychological Science}, 11(4):512--534, 2016.

\bibitem[Mellers et al.(2014)]{mellers2014}
Barbara Mellers, Lyle Ungar, Jonathan Baron, Jaime Ramos, Burcu Gurcay, Katrina Fincher, et al.
\newblock Psychological strategies for winning a geopolitical forecasting tournament.
\newblock \emph{Psychological Science}, 25(5):1106--1115, 2014.

\bibitem[Miall et al.(2007)]{miall2007}
R. Chris Miall, Lars O. D. Christensen, Owen Cain, and James Stanley.
\newblock Disruption of state estimation in the human lateral cerebellum.
\newblock \emph{PLoS Biology}, 5(11):e316, 2007.

\bibitem[Miall and Wolpert(1996)]{miallwolpert1996}
R. Chris Miall and Daniel M. Wolpert.
\newblock Forward models for physiological motor control.
\newblock \emph{Neural Networks}, 9(8):1265--1279, 1996.

\bibitem[Miloyan et al.(2019)]{miloyan2019}
Beyon Miloyan, Kimberley A. McFarlane, and Thomas Suddendorf.
\newblock Measuring mental time travel: is the hippocampus really critical for episodic memory and episodic foresight?
\newblock \emph{Cortex}, 117:371--384, 2019.

\bibitem[Moeck et al.(2026)]{moeck2026}
Ella K. Moeck, Karta K. Grewal, Amit Mehta, Katharine H. Greenaway, Peter Koval, and Elise K. Kalokerinos.
\newblock Affective forecasting accuracy in everyday life.
\newblock \emph{Affective Science}, 7:306--318, 2026.

\bibitem[Momennejad et al.(2017)]{momennejad2017}
Ida Momennejad, Evan M. Russek, Jin H. Cheong, Matthew M. Botvinick, Nathaniel D. Daw, and Samuel J. Gershman.
\newblock The successor representation in human reinforcement learning.
\newblock \emph{Nature Human Behaviour}, 1(9):680--692, 2017.

\bibitem[M\"uller et al.(2006)]{muller2006}
Sean M\"uller, Bruce Abernethy, and Damian Farrow.
\newblock How do world-class cricket batsmen anticipate a bowler's intention?
\newblock \emph{Quarterly Journal of Experimental Psychology}, 59(12):2162--2186, 2006.

\bibitem[M\"uller et al.(2010)]{muller2010}
Sean M\"uller, Bruce Abernethy, Jason Reece, Melissa Rose, et al.
\newblock An in-situ examination of the timing of information pick-up for interception by cricket batsmen of different skill levels.
\newblock \emph{Perception}, 39(6):745--760, 2010.

\bibitem[Rachlin and Green(1972)]{rachlin1972}
Howard Rachlin and Leonard Green.
\newblock Commitment, choice and self-control.
\newblock \emph{Journal of the Experimental Analysis of Behavior}, 17(1):15--22, 1972.

\bibitem[Rakitin et al.(1998)]{rakitin1998}
Brian C. Rakitin, John Gibbon, Trevor B. Penney, Chara Malapani, Sean C. Hinton, and Warren H. Meck.
\newblock Scalar expectancy theory and peak-interval timing in humans.
\newblock \emph{Journal of Experimental Psychology: Animal Behavior Processes}, 24(1):15--33, 1998.

\bibitem[Roy(1952)]{roy1952}
A. D. Roy.
\newblock Safety first and the holding of assets.
\newblock \emph{Econometrica}, 20(3):431--449, 1952.

\bibitem[Rozenblit and Keil(2002)]{rozenblit2002}
Leonid Rozenblit and Frank Keil.
\newblock The misunderstood limits of folk science: an illusion of explanatory depth.
\newblock \emph{Cognitive Science}, 26(5):521--562, 2002.

\bibitem[Sala and Gobet(2017a)]{sala2017cdps}
Giovanni Sala and Fernand Gobet.
\newblock Does far transfer exist? Negative evidence from chess, music, and working memory training.
\newblock \emph{Current Directions in Psychological Science}, 26(6):515--520, 2017.

\bibitem[Sala and Gobet(2017b)]{sala2017lb}
Giovanni Sala and Fernand Gobet.
\newblock Working memory training in typically developing children: a meta-analysis of the available evidence.
\newblock \emph{Learning \& Behavior}, 45(4):414--421, 2017.

\bibitem[Sala and Gobet(2017c)]{sala2017mc}
Giovanni Sala and Fernand Gobet.
\newblock Experts' memory superiority for domain-specific random material generalizes across fields of expertise.
\newblock \emph{Memory \& Cognition}, 45(2):183--193, 2017.

\bibitem[Scholz and Sch\"oner(1999)]{scholz1999}
John P. Scholz and Gregor Sch\"oner.
\newblock The uncontrolled manifold concept: identifying control variables for a functional task.
\newblock \emph{Experimental Brain Research}, 126(3):289--306, 1999.

\bibitem[Selten(1998)]{selten1998}
Reinhard Selten.
\newblock Aspiration adaptation theory.
\newblock \emph{Journal of Mathematical Psychology}, 42(2--3):191--214, 1998.

\bibitem[Simon(1956)]{simon1956}
Herbert A. Simon.
\newblock Rational choice and the structure of the environment.
\newblock \emph{Psychological Review}, 63(2):129--138, 1956.

\bibitem[Simons and Chabris(1999)]{simons1999}
Daniel J. Simons and Christopher F. Chabris.
\newblock Gorillas in our midst: sustained inattentional blindness for dynamic events.
\newblock \emph{Perception}, 28(9):1059--1074, 1999.

\bibitem[Snider et al.(2015)]{snider2015}
Joseph Snider, Dongpyo Lee, Howard Poizner, and Sergei Gepshtein.
\newblock Prospective optimization with limited resources.
\newblock \emph{PLoS Computational Biology}, 11(9):e1004501, 2015.

\bibitem[Soderberg et al.(2015)]{soderberg2015}
Courtney K. Soderberg, Shannon P. Callahan, Annie O. Kochersberger, Elinor Amit, and Alison Ledgerwood.
\newblock The effects of psychological distance on abstraction: two meta-analyses.
\newblock \emph{Psychological Bulletin}, 141(3):525--548, 2015.

\bibitem[Squire et al.(2010)]{squire2010}
Larry R. Squire, Anna S. van der Horst, Susan G. R. McDuff, Jennifer C. Frascino, Ramona O. Hopkins, and Kristin N. Mauldin.
\newblock Role of the hippocampus in remembering the past and imagining the future.
\newblock \emph{Proceedings of the National Academy of Sciences}, 107(44):19044--19048, 2010.

\bibitem[St\"uttgen et al.(2012)]{stuttgen2012}
Peter St\"uttgen, Peter Boatwright, and Robert T. Monroe.
\newblock A satisficing choice model.
\newblock \emph{Marketing Science}, 31(6):878--899, 2012.

\bibitem[Todorov and Jordan(2002)]{todorov2002}
Emanuel Todorov and Michael I. Jordan.
\newblock Optimal feedback control as a theory of motor coordination.
\newblock \emph{Nature Neuroscience}, 5(11):1226--1235, 2002.

\bibitem[Trope and Liberman(2010)]{tropeliberman2010}
Yaacov Trope and Nira Liberman.
\newblock Construal-level theory of psychological distance.
\newblock \emph{Psychological Review}, 117(2):440--463, 2010.

\bibitem[Warren(1984)]{warren1984}
William H. Warren.
\newblock Perceiving affordances: visual guidance of stair climbing.
\newblock \emph{Journal of Experimental Psychology: Human Perception and Performance}, 10(5):683--703, 1984.

\bibitem[Warren and Whang(1987)]{warrenwhang1987}
William H. Warren and Suzanne Whang.
\newblock Visual guidance of walking through apertures: body-scaled information for affordances.
\newblock \emph{Journal of Experimental Psychology: Human Perception and Performance}, 13(3):371--383, 1987.

\bibitem[Wilde(1982)]{wilde1982}
Gerald J. S. Wilde.
\newblock The theory of risk homeostasis: implications for safety and health.
\newblock \emph{Risk Analysis}, 2(4):209--225, 1982.

\bibitem[Wilson and Gilbert(2005)]{wilsongilbert2005}
Timothy D. Wilson and Daniel T. Gilbert.
\newblock Affective forecasting: knowing what to want.
\newblock \emph{Current Directions in Psychological Science}, 14(3):131--134, 2005.

\bibitem[Zeki and Bal\c{c}\i(2019)]{zeki2019}
Mustafa Zeki and Fuat Bal\c{c}\i.
\newblock A simplified model of communication between time cells.
\newblock \emph{Frontiers in Computational Neuroscience}, 12:111, 2019.
\end{thebibliography}

\end{document}
